% --- Archon physics formalization source begin ---
% archon:physics
% archon:chemistry
% archon:covers IChO2026Problems/problem_icho_2026_t5_a4.lean
% archon:source-report reports/icho_2026/problem_icho_2026_t5_a4.source.json
% archon:problem-id icho_2026_t5
% archon:part-id T5-A4
% archon:previous-part icho_2026_t5_a3
% archon:previous-part-policy natural_language_prerequisite_only

\chapter{Icho Chemistry problem icho\_2026\_t5\_a4}
\label{ch:IChO2026Problems_problem_icho_2026_t5_a4}

\paragraph{Problem source.}
\#\# Source
58th International Chemistry Olympiad, Tashkent, Uzbekistan, 2026.
Official-materials index: https://www.icho2026.uz/problems
English problem PDF: ../icho\_2026\_source/raw/theory\_problem.pdf
English solution/rubric PDF: ../icho\_2026\_source/raw/theory\_solution.pdf
Problem PDF page 46; printed local question page 3; solution starts on PDF page 43.
The page PNG is primary visual evidence for formulas, structures, tables, and figures.

\#\# T5. Cardiolipins
T5. Cardiolipins 6\% In this problem you are asked to be creative and build molecules from fragments. For instance, if you are required to draw a chiral phospholipid such as W using fragments a-e, it can be assembled as shown below. Cardiolipins are a family of acyclic phospholipids differing in fatty acid residues. They are essential for energy production and membrane stability. Defects can lead to certain diseases. PL1 belongs to this family. Using only the structural elements a-d in the quantities stated below, it is possible to assemble the non-ionised form of PL1. R is a hydrocarbon substituent in the fatty acid structure. PL1 does not contain any peroxide bonds. 5.1 Tick one correct statement regarding the value of n (number of type a fragments). (a) n is an even number, (b) n is an odd number, (c) n can be either an odd or an even number. Some cardiolipins are chiral molecules, even if they contain four identical fatty acid residues, as in the case of PL1. One enantiomer of PL1 is found in prokaryotes and eukaryotes, and the other only in archaea. All other diastereomers of PL1 are achiral molecules, since they have a plane of symmetry. Do not consider the chirality of phosphorus atoms as stereocentres. PL1 is a diprotic acid with the same acidic groups. Its second deprotonation is less favourable (𝑝𝐾𝑎2 ≫𝑝𝐾𝑎1) due to the special, stable structure of monoanion Y. 5.2 Draw the structure of one enantiomer of PL1. Draw the structure of Y in a way that shows the stabilisation that explains the difference between 𝑝𝐾𝑎1 and 𝑝𝐾𝑎2. Use the abbreviation R for the fatty acid residues. A series of experiments was conducted to identify the fatty acid (RCOOH) in PL1 found predominantly in mammalian hearts. During reductive ozonolysis, RCOOH forms three different organic products in equimolar amounts. 5.3 Determine the molecular formula of the fatty acid (RCOOH), if the non-ionised form of PL1 contains 255 𝜎and 𝜋bonds between atoms in total. Support your answer with calculations. If you were unable to find the structural formula of PL1, you can use a-d fragments from 5.1. 100 g of RCOOH reacts with 181.0 g of iodine. RCOOH also reacts in similar way with X and forms an adduct with an iodine mass fraction of 36.57\%.

\#\# Current subquestion 5.4
Determine the molecular formula of X. Support your answer with calculations.

\#\# Dataset metadata
IChO 2026; T5; raw rubric points=2; kind=theory; formalization\_ready=true.

\paragraph{Shared problem context.}
T5. Cardiolipins 6\% In this problem you are asked to be creative and build molecules from fragments. For instance, if you are required to draw a chiral phospholipid such as W using fragments a-e, it can be assembled as shown below. Cardiolipins are a family of acyclic phospholipids differing in fatty acid residues. They are essential for energy production and membrane stability. Defects can lead to certain diseases. PL1 belongs to this family. Using only the structural elements a-d in the quantities stated below, it is possible to assemble the non-ionised form of PL1. R is a hydrocarbon substituent in the fatty acid structure. PL1 does not contain any peroxide bonds. 5.1 Tick one correct statement regarding the value of n (number of type a fragments). (a) n is an even number, (b) n is an odd number, (c) n can be either an odd or an even number. Some cardiolipins are chiral molecules, even if they contain four identical fatty acid residues, as in the case of PL1. One enantiomer of PL1 is found in prokaryotes and eukaryotes, and the other only in archaea. All other diastereomers of PL1 are achiral molecules, since they have a plane of symmetry. Do not consider the chirality of phosphorus atoms as stereocentres. PL1 is a diprotic acid with the same acidic groups. Its second deprotonation is less favourable (𝑝𝐾𝑎2 ≫𝑝𝐾𝑎1) due to the special, stable structure of monoanion Y. 5.2 Draw the structure of one enantiomer of PL1. Draw the structure of Y in a way that shows the stabilisation that explains the difference between 𝑝𝐾𝑎1 and 𝑝𝐾𝑎2. Use the abbreviation R for the fatty acid residues. A series of experiments was conducted to identify the fatty acid (RCOOH) in PL1 found predominantly in mammalian hearts. During reductive ozonolysis, RCOOH forms three different organic products in equimolar amounts. 5.3 Determine the molecular formula of the fatty acid (RCOOH), if the non-ionised form of PL1 contains 255 𝜎and 𝜋bonds between atoms in total. Support your answer with calculations. If you were unable to find the structural formula of PL1, you can use a-d fragments from 5.1. 100 g of RCOOH reacts with 181.0 g of iodine. RCOOH also reacts in similar way with X and forms an adduct with an iodine mass fraction of 36.57\%.

\paragraph{Current subquestion.}
Determine the molecular formula of X. Support your answer with calculations.

\paragraph{Recorded answer/context.}
During determination of the iodine number, the following reaction takes place: The iodine number of RCOOH and the iodine mass fraction in the adduct are described as follows: 181.0 = 𝑁C=C×253.8 𝑀RCOOH × 100 ⇒𝑀RCOOH 𝑁C=C = 140.2 0.3657 = 𝑁C=C×126.9 𝑀RCOOH+𝑁C=C×(126.9+𝐴Hal) ⇒𝐴Hal = 79.9 g mol−1. So, X is iodine monobromide, IBr. If X was molecular iodine (I2), then: 𝜔𝐼= 𝑁C=C×253.8 𝑀RCOOH+𝑁C=C×253.8 = 0.6441 which is incorrect. For correct formula of X: 1 point. For correct calculations: 1 point. 2.0 pt

\paragraph{Figure/image paths.}
\begin{itemize}
\item ../icho\_2026\_source/image/T5\_page-3.png
\item ../icho\_2026\_source/image/T5\_page-2.png
\end{itemize}

\paragraph{Reusable previous-part conclusions.}
\begin{itemize}
\item Source T5-A3. Question: Determine the molecular formula of the fatty acid (RCOOH), if the non-ionised form of PL1 contains 255 𝜎and 𝜋bonds between atoms in total. Support your answer with calculations. If you were unable to find the structural formula of PL1, you can use a-d fragments from 5.1. Reusable conclusions: Correct scheme of reductive ozonolysis of RCOOH is as follows: So, 𝑁𝐶=𝐶= 2. If we take R as 𝐶𝑥𝐻𝑦, then: y=(2x+1)-𝑁𝐶=𝐶∙2=2x-3.(eq.3.1) Then the formula for calculating the number of bonds in PL1 can be written as follows: 0.5∙1+5∙2+11.5∙3+4∙(3+2x+0.5y)=255, 4x+y=99.(eq.3.2) Combining eq.3.1 and eq.3.2: 4x+y=4x+2x-3=6x-3=99 ⟹x=17, y=31. RCOOH −C17H31COOH or C18H32O2 For correct eq.3.2: 2 points. For correct number of double bonds: 1 point. For correct formula of C17H31COOH: 1 point. 4.0 pt Policy: natural\_language\_prerequisite\_only; the current target and later answers must not be assumed
\end{itemize}

\paragraph{Formalization target.}
The verified Lean formalization is in `IChO2026Problems/problem\_icho\_2026\_t5\_a4.lean`,
with its theorem and lemma proofs complete.
The Lean declarations must preserve every mathematical object, quantifier, hypothesis, side condition, approximation/error guarantee, and requested conclusion from the source statement.
The implementation reuses the pinned Mathlib, Physlib, and Chemistry libraries,
with local definitions only for source-specific objects.

\begin{definition}[Iodine atomic mass, candidate halogens, and interhalogen formulas]
\label{def:physics:icho_2026_t5_a4:setup}
\lean{IChO2026T5A4.ArI, IChO2026T5A4.Halogen, IChO2026T5A4.Halogen.atomicMass, IChO2026T5A4.MolecularFormula, IChO2026T5A4.MolecularFormula.IHal, IChO2026T5A4.MolecularFormula.IBr, IChO2026T5A4.MolecularFormula.molarMass, IChO2026T5A4.MolecularFormula.iodineMass, IChO2026T5A4.ReagentX}
\leanok
The standard atomic mass of iodine is $\mathrm{ArI} = 126.9$ g·mol${}^{-1}$
(exam periodic table), so the molar mass of $\mathrm{I_2}$ is
$2\,\mathrm{ArI} = 253.8$ g·mol${}^{-1}$.  The unknown reagent $X$ reacts
with RCOOH "in a similar way" to iodine, so it is modeled as a diatomic
interhalogen $\mathrm{I\!-\!Hal}$ whose second atom ranges over the four
candidate halogens F, Cl, Br, I (the I branch being molecular iodine
$\mathrm{I_2}$ itself).  A molecular formula is an explicit atom count per
element; its molar-mass law sums count times standard atomic mass, and its
iodine-content law counts the iodine atoms times $\mathrm{ArI}$.
\end{definition}

\begin{lemma}[Mass laws of an interhalogen $\mathrm{I\!-\!Hal}$]
\label{lem:physics:icho_2026_t5_a4:masslaws}
\lean{IChO2026T5A4.MolecularFormula.molarMass_IHal, IChO2026T5A4.MolecularFormula.iodineMass_IHal}
\leanok
For every candidate halogen $h$, the formula $\mathrm{I\!-\!}h$ has molar
mass $\mathrm{ArI} + \mathrm{Ar}(h)$ (in particular
$M(\mathrm{I_2}) = 2\,\mathrm{ArI}$) and carries
$1$ mole of iodine atoms per mole ($2$ for $\mathrm{I_2}$).
\end{lemma}
\begin{proof}
\leanok
Case analysis on $h$: the atom counts of $\mathrm{I\!-\!}h$ evaluate by
constructor computation, and both sides reduce to the same sum of tabulated
atomic masses.
\end{proof}

\begin{theorem}[Iodine-number experiment: molar mass per double bond]
\label{thm:physics:icho_2026_t5_a4:iodinenumber}
\lean{IChO2026T5A4.molarMass_of_iodine_experiment}
\leanok
If $100$ g of an acid of molar mass $M$ with $N = 2$ C=C bonds consumes
$181.0$ g of iodine at one molecule of $\mathrm{I_2}$ per double bond, i.e.\
$N \cdot (100/M) \cdot 2\,\mathrm{ArI} = 181.0$, then
$M = N \cdot 2\,\mathrm{ArI} \cdot 100 / 181.0$ and
$|M/N - 140.2| \le 0.03$ (exact value $M/N = 25380/181 = 140.221\ldots$).
\end{theorem}
\begin{proof}
\leanok
Clearing the denominator $M > 0$ in the iodine-number relation gives
$181 \cdot M = 2 \cdot 100 \cdot 2\,\mathrm{ArI} = 50760$, so
$M = 50760/181$ and $M/N = 25380/181 = 140.2209\ldots$, within $0.03$ of the
marking scheme's rounded $140.2$.
\end{proof}

\begin{theorem}[Experimental atomic mass of the second halogen]
\label{thm:physics:icho_2026_t5_a4:halogenmass}
\lean{IChO2026T5A4.experimental_halogen_mass_exact, IChO2026T5A4.experimental_halogen_mass_approx}
\leanok
If the adduct's iodine mass fraction obeys
$N \cdot \mathrm{ArI} / (M + N \cdot (\mathrm{ArI} + A)) = 0.3657$ together
with the iodine-number relation, then the second halogen's experimental
atomic mass is
$A = \mathrm{ArI}/0.3657 - \mathrm{ArI} - 2\,\mathrm{ArI}\cdot 100/181.0
= 79.885\ldots$ g·mol${}^{-1}$, i.e.\ $|A - 79.9| \le 0.02$.
\end{theorem}
\begin{proof}
\leanok
The denominator $M + N(\mathrm{ArI} + A)$ is positive, so the mass-fraction
equation rearranges to $A = \mathrm{ArI}/0.3657 - \mathrm{ArI} - M/N$;
substituting $M/N = 25380/181$ from the iodine-number experiment gives the
closed form, and exact rational evaluation yields
$A = 79.8847\ldots$, within $0.02$ of the printed $79.9$.
\end{proof}

\begin{theorem}[Requested conclusion: $X$ is iodine monobromide $\mathrm{IBr}$]
\label{thm:physics:icho_2026_t5_a4:target}
\lean{IChO2026T5A4.X_formula_eq_IBr}
\leanok
If $X = \mathrm{I\!-\!}h$ for some candidate halogen $h$ and the
formula-derived iodine mass fraction of the adduct matches the measured
(rounded) value $0.3657$ within half a unit of the last digit
($|\omega - 0.3657| \le 0.00005$), then $X$ is $\mathrm{IBr}$.
\end{theorem}
\begin{proof}
\leanok
Case analysis on $h$.  With $M = 50760/181$ fixed by the iodine-number
experiment, the adduct fractions for $\mathrm{IF}$, $\mathrm{ICl}$ and
$\mathrm{I_2}$ evaluate exactly to $0.4435\ldots$, $0.4194\ldots$ and
$181/281 = 0.6441\ldots$, each violating the $0.00005$ tolerance, while
$\mathrm{IBr}$ gives $0.36568\ldots$, within tolerance; definitionally
$\mathrm{IBr} = \mathrm{I\!-\!Br}$.
\end{proof}

\begin{theorem}[Consistency check: the $\mathrm{I_2}$ branch is refuted]
\label{thm:physics:icho_2026_t5_a4:diiodine}
\lean{IChO2026T5A4.diiodine_adduct_fraction}
\leanok
Had $X$ been molecular iodine $\mathrm{I_2}$, the adduct's iodine mass
fraction would be exactly
$N \cdot 2\,\mathrm{ArI} / (M + N \cdot 2\,\mathrm{ArI}) = 181/281
= 0.6441\ldots \neq 0.3657$.
\end{theorem}
\begin{proof}
\leanok
Substituting $M = 50760/181$ and $\mathrm{ArI} = 126.9$ makes the fraction a
concrete rational: $507.6 \cdot 181 / (50760 + 507.6 \cdot 181) = 181/281$,
which differs from $0.3657$.
\end{proof}
% --- Archon physics formalization source end ---
